\reportheading{6.\quad 结论与设计推荐}
仅开启一扇门时，交换流集中在 R1，实测顶棚温度峰值向封闭侧偏移。用 15~mm 导热板替代绝热壁后，106.2~kW 的 OP-6 顶棚温度剖面更接近实验值。47.2~kW 的 OP-1 峰值位置保持在燃烧器上方，预测峰值温升降低。

对于第 4 节验证的丙烯火排烟布局，推荐总排烟量为 1.12~m$^3$/s。六组细网格 FDS 计算中，八个测点在 120~s 内均满足 0.38~m 净高要求，最小余量为 55.8~mm。在三个已测风量中，该方案满足净高约束，且所需的支管估算电功率最低。

BO 和 GA 在原布局的 16 个方案中比固定顺序更早找到净高余量最大的组合。在改进布局的细网格归档数据上，固定风量次序、仅用 PINN 排工况需 7--11 次检查。由 PINN 同时排风量和工况时，确认最低功率可行风量需 9--17 次，固定基线为 7 次。用 BO 或 GA 与代理模型联动调度新 FDS 计算，仍有待实际运行检验。
